%% notes-tex/interaction-fitness-figure-gate/sections/4-remaining.tex
%% SOURCE: hand-authored; the values named are macros of tables/gate_numbers.tex.

\section{What is established and what remains\secstatus{todo}}
\label{sec:remaining}

\pfstep{Established} The published trigenic score is the identity of
equation~\ref{eq:identity} to the rounding floor on both screens, so a held-out triple's
null expectation is a function of training records (\file{025-s3-closure}). The free
arithmetic is measured on both arms and is large: \gloRNOnemultquerydoubletest{} and
\gloQNOnemultquerydoubletest{} with no parameters, \gloRLTwohgblowerordertest{} and
\gloQLTwohgblowerordertest{} with boosted trees. The transformer as trained reads below
the zero-parameter null on both arms (\gcgtRValWindow{}, \gcgtQValWindow{}, validation,
one seed). The gene-identity ridge sits at \gloRBOneadditivegenetest{} and
\gloQBOneadditivegenetest{}, so the transformer does beat the ladder the earlier draft
compared it to, and that comparison is the wrong one. The fitness of a held-out triple is
predictable from the lower-order terms at \gloRFOneridgelowerorderfitnesstest{} on R.

\pfstep{Not established} Whether a transformer that reads the lower-order terms as inputs
beats the boosted trees, which is the only comparison in which the two regimes match
(E1). Whether the transformer's residual over the digenic-controls null carries
anything, read as a target (E2). Any test score of a 030 checkpoint, and any seed
spread; every transformer number here is one seed on validation (E3, E4). What the
ladder reads at deployment tiers, where a new triple has its singles but not its
doubles (E5). The comparators on the same data (E7). Why the arithmetic explains far
less on the Kuzmin 2018 rows than on the 2020 rows
(\gloRNOnemultquerydoubletestKzTwoZeroOneEight{} against
\gloRNOnemultquerydoubletestKzTwoZeroTwoZero{} on arm R test); hypothesis, untested:
the 2020 diagnostic-array screen measured its control queries on the same strains in the
same batch as the triples, so the control scores carry the batch, and the 2018
genome-wide screen did not.

\pfstep{Corrections recorded} The 025 document's decision table compared the transformer
to the additive ridge and found it below that ridge on arm Q test for one checkpoint; that
reading stands, and this document adds that the ridge itself is far under the lower-order
null. The weekly note of 2026-10-11 read the Q transformer at 0.391 over epochs 19 to 28;
the registry has since synced to epoch \gcgtQLastEpoch{} and the window moved to
\gcgtQValWindow{}; both are one seed, validation.
